\section{Конечные алгебры и машинная арифметика. Двоичный смещенный код. «Арифметический круг». Двоичная и логические арифметики. Примеры.}

\subsection*{Конечные алгебры в ЭВМ}
Машинная арифметика — это реализация кольца $(\mathbb{Z}_n; +, \ast)$, где $n = 2^k$, а $k$ — разрядность процессора (16, 32, 64 бита). В отличие от бесконечных целых чисел $\mathbb{Z}$, машинная арифметика конечна и имеет специфические свойства (наличие делителей нуля, переполнение).

Для визуализации операций в $\mathbb{Z}_{2^k}$ используется \textbf{арифметический круг}. Элементы располагаются по кругу от $0$ до $2^k-1$. 
\begin{itemize}
    \item Сложение соответствует движению по часовой стрелке.
    \item Вычитание — против часовой стрелки.
    \item \textbf{Переполнение} возникает, когда результат операции «перескакивает» через границу $2^k-1 \to 0$.
\end{itemize}
В знаковой арифметике круг делится пополам: первая половина $[0, 2^{k-1}-1]$ — положительные числа, вторая половина $[2^{k-1}, 2^k-1]$ — отрицательные числа (в дополнительном коде).

\paragraph{Двоичный смещенный код}
Для представления вещественных чисел (экспоненты) или упрощения сравнения целых используется \textbf{смещенный код}. 
Число $x$ представляется как $x_{code} = x + B$, где $B$ — фиксированное смещение (обычно $B = 2^{k-1}$).
\begin{itemize}
    \item Это позволяет представлять отрицательные числа натуральными, сохраняя порядок (большему числу соответствует больший код).
\end{itemize}

\paragraph{Двоичная и логические арифметики}
Новиков разделяет операции над машинными словами на два типа:
\begin{enumerate}
    \item \textbf{Двоичная арифметика:} Операции в кольце $(\mathbb{Z}_{2^k}; +, \ast)$. Здесь учитываются переносы между разрядами.
    \item \textbf{Логические арифметики:} Поразрядные операции ($\land, \lor, \oplus, \neg$). В этой структуре каждый бит независим, и перенос не возникает. Это алгебра $(\mathbb{B}^k; \land, \lor, \neg)$.
\end{enumerate}

\paragraph{Пример}
В 8-битной арифметике ($Z_{256}$):
$200 + 100 = 300 \pmod{256} = 44$. (Двоичная арифметика).
$11001000_2 \text{ AND } 01100100_2 = 01000000_2$. (Логическая арифметика).